% Minimal retained preamble for lane-p2 cards taken from
% sources/logic-fl02-conservation-of-ramsey-s-theorem-f/pinned-original-source.tex
% (Conservation of Ramsey's theorem for pairs, arXiv:2402.11616v3).
%
% Copied verbatim from the author's own file:
%   line 2      : \documentclass{article}
%   lines 16-25 : the author's theorem declarations
%   lines 46-88 : the author's macro definitions for the reverse-mathematics principles
%   line 91     : \card
%   lines 104-107: the natural sums and products, which rely on mathabx (source line 103)
%   line 116    : \qt
%
% Omitted author preamble lines and why (nothing was renamed or redefined):
%   line 3 (graphicx), 7 (url), 9 (cancel): no figure, url or cancel use in the retained spans;
%   line 6 (cleveref): the assembler adds its own cleveref machinery, a second load clashes;
%   line 8 (\usepackage[utf8x]{inputenc}): inputenc with the utf8x option fights the forced
%       xelatex engine;
%   lines 28-37 (the rep@theorem machinery and \newreptheorem{maintheorem}): the retained spans
%       use no rep-theorem environment;
%   lines 40-43 (xcolor and the three author comment macros): not used by the retained spans;
%   lines 93-99 (the LS2 font encoding block for \smashtimes based on stix): the host TeX Live has
%       no stix directory and the author himself never uses the result -- line 105 maps \ntimes to
%       \dottimes, and line 110 shows \smashtimes already commented out in the source.
% mathabx is retained because \dotplus, \dottimes and \dotdiv come from it.
\documentclass{article}
\usepackage{amsmath,amssymb,amsthm}
\usepackage{mathabx}
\usepackage{cancel}
\usepackage{cleveref}

\newtheorem{theorem}{Theorem}
\numberwithin{theorem}{section}
\newtheorem{maintheorem}[theorem]{Main Theorem}
\newtheorem{lemma}[theorem]{Lemma}
\newtheorem{question}[theorem]{Question}
\newtheorem{proposition}[theorem]{Proposition}
\newtheorem{statement}[theorem]{Statement}
\newtheorem{remark}[theorem]{Remark}
\newtheorem{definition}[theorem]{Definition}
\newtheorem{corollary}[theorem]{Corollary}

\newcommand{\RCA}{\mathsf{RCA}}
\newcommand{\WKL}{\mathsf{WKL}}
\newcommand{\SRT}{\mathsf{SRT}}
\newcommand{\EM}{\mathsf{EM}}
\newcommand{\SEM}{\mathsf{SEM}}
\newcommand{\ADS}{\mathsf{ADS}}
\newcommand{\SCAC}{\mathsf{SCAC}}
\newcommand{\SADS}{\mathsf{SADS}}
\newcommand{\COH}{\mathsf{COH}}
\newcommand{\ACA}{\mathsf{ACA}}
\newcommand{\qvdash}{\operatorname{{?}{\vdash}}}
\newcommand{\nqvdash}{\operatorname{{?}{\nvdash}}}
\newcommand{\BSig}{\mathsf{B}\Sigma^0}
\newcommand{\ISig}{\mathsf{I}\Sigma^0}
\newcommand{\LSig}{\mathsf{L}\Sigma^0}
\newcommand{\BPi}{\mathsf{B}\Pi^0}
\newcommand{\IPi}{\mathsf{I}\Pi^0}
\newcommand{\LPi}{\mathsf{L}\Pi^0}
\newcommand{\IDel}{\mathsf{I}\Delta^0}
\newcommand{\BCDe}{\mathsf{BC}\Delta^0}
\newcommand{\WF}{\mathsf{WF}}
\newcommand{\RT}{\mathsf{RT}}
\newcommand{\BME}{\mathsf{BME}}
\newcommand{\TT}{\mathsf{TT}}
\renewcommand{\O}{\mathcal{O}}
\renewcommand{\L}{\mathcal{L}}
\newcommand{\C}{\mathcal{C}}
\newcommand{\R}{\mathcal{R}}
\renewcommand{\S}{\mathcal{S}}
\newcommand{\T}{\mathcal{T}}
\renewcommand{\P}{\mathcal{P}}
\newcommand{\M}{\mathcal{M}}
\newcommand{\Nc}{\mathcal{N}}
\renewcommand{\P}{\mathcal{P}}
\newcommand{\Q}{\mathcal{Q}}
\newcommand{\NN}{\mathbb{N}}
\newcommand{\PP}{\mathbb{P}}
\newcommand{\uh}[0]{{\upharpoonright}}
\newcommand{\card}{\operatorname{card}}
\newcommand{\nplus}{\dotplus}
\newcommand{\ntimes}{\dottimes}
\newcommand{\nminus}{\dotdiv}
\newcommand{\nsum}{\dot{\sum}}
\def\qt#1{``#1''}
\title{A finite maximal tree of split pairs forces a last true instance on one of the two sides}
\author{Quentin Le Houérou, Ludovic Levy Patey, Keita Yokoyama}
\date{arxiv:2402.11616v3; CC BY 4.0}
\begin{document}
\maketitle
\noindent\textit{Editorial note.} The counted unit is the article's own Lemma lem:sads-wf-finite-tree-forcing at source lines 1206-1215 together with the proof environment that begins on the next line (1216-1231); the pairing is the ordinary adjacent one. The author has left a commented-out earlier version of the same lemma immediately after the proof at source lines 1234 onwards; it is not author text for our purposes and is not retained.\par\medskip
\noindent\textit{Editorial note.} The retained context at source lines 1149-1156 is the statement of the forcing question lemma, which the proof invokes twice, and it is counted by a sibling card of this lane; the retained context at source lines 1169-1183 and 1186-1192 is the author's own definitions of trees of split pairs, validity and satisfaction.\par\medskip
\noindent\textit{Editorial note.} The retained preamble reproduces verbatim the author's document class at source line 2, his theorem declarations at source lines 16-25, his macro definitions at source lines 46-88 and 91, and his natural-operation macros at source lines 104-107, and it retains mathabx and cleveref; the stix-dependent LS2 block and the inputenc load with the utf8x option are not retained.\par\medskip
\noindent\textit{Editorial note.} No retained line contains a citation, so the card carries no bibliography; the only cross-references of the retained proof are the two author Cref calls to the forcing question lemma, whose statement is retained in the context.\par\medskip
\noindent\textit{Editorial note.} The article's conservation theorems, its COH, WKL and SEM lemmas and its other SADS lemmas are not counted here.\par\medskip
\noindent\textit{Allowed standard prerequisites (editorial).}\begin{itemize}
\item Trees of split pairs, validity and satisfaction. A tree attached to a condition is a binary tree labelled by split pairs extending the condition, with successors extending the two components; a node is valid when along the left branches above it the ascending components lie inside U and along the right branches the descending components lie inside V; a node satisfies a pair of Sigma\textasciicircum{}0\_1 formulas when the first holds of the ascending component at the node rank and the second of the descending component at the number of ones, and a tree satisfies the pair when all its nodes do; the author's own definitions are retained at source lines 1169-1183 and 1186-1192.. Fixes the objects and the satisfaction property of the counted lemma, and supplies the validity notion used to pick the node whose failure of an extension contradicts maximality.
\item The forcing question and its two-case decision. For a condition and two computations, the forcing question relates the two sides through a split pair, and the article proves that if the question fails then some extension forces divergence on one of the two sides, decidably so; retained with this card at source lines 1149-1156.. The proof uses it twice: once when the tree is empty, and once at the maximal valid node, to convert the failure of the forcing question into an extension forcing the negation of the next instance.
\item Induction along the tree by L Delta\textasciicircum{}0\_2. The scheme of induction for L Delta\textasciicircum{}0\_2 formulas, available because the model satisfies collection for Sigma\textasciicircum{}0\_2, allows a maximal valid element of a finite tree to be taken.. The proof selects a valid node with no valid descendant, which needs that scheme.
\end{itemize}\medskip
\begin{lemma}\label[lemma]{lem:sads-forcing-question-sigma01}
    Let $(\sigma_0, \tau_0)$ be a condition, and let $e_0, e_1, x_0, x_1 \in M$. Then:
    \begin{itemize}
        \item If $(\sigma_0, \tau_0) \qvdash \Phi_{e_0}^{G_0 \oplus Y}(x_0) \downarrow \vee \Phi_{e_1}^{G_1 \oplus Y}(x_1) \downarrow$, then there exist an extension $(\sigma_1, \tau_1) \leq (\sigma_0, \tau_0)$ and some $i < 2$ such that $(\sigma_1, \tau_1) \Vdash_i \Phi_{e_i}^{G_i \oplus Y}(x_i) \downarrow$.
        \item If $(\sigma_0, \tau_0) \nqvdash \Phi_{e_0}^{G_0 \oplus Y}(x_0) \downarrow \vee \Phi_{e_1}^{G_1 \oplus Y}(x_1) \downarrow$, then there exist an extension $(\sigma_1, \tau_1) \leq (\sigma_0, \tau_0)$ and some $i < 2$ such that $(\sigma_1, \tau_1) \Vdash_i \Phi_{e_i}^{G_i \oplus Y}(x_i) \uparrow$.
    \end{itemize}
    Furthermore, deciding which case holds and computing the extension $(\sigma_1, \tau_1)$ can be done $Y'$-computably.
\end{lemma}

\begin{definition}
Fix a condition $(\sigma, \tau)$.

A \emph{$(\sigma, \tau)$-tree} is a binary tree $T$ labelled by a family of split pairs $(\sigma_{\rho}, \tau_{\rho})_{\rho \in T}$ such that for all $\rho \in T$:
\begin{itemize}
    \item $\sigma \preceq \sigma_{\rho}$ and $\tau \preceq \tau_{\rho}$ ;
    \item For every $\mu \in T$, if $\rho \cdot 0 \preceq \mu$ then $\sigma_{\rho} \preceq \sigma_{\mu}$ and if $\rho \cdot 1 \preceq \mu$ then $\tau_{\rho} \preceq \tau_{\mu}$ ;
\end{itemize}
%    For a node $\rho \in T$, we define $\pi(\rho)$ to be equal to $(\sigma_{\rho \uh i_0}, \tau_{\rho \uh j_0})$ where $i_0$ is the biggest index $i < |\rho|$ such that $\rho(i) = 0$ (if no such index exists then we take $\sigma$ instead of $\sigma_{\rho \uh i_0}$ in the definition of $\pi(\rho)$) and $j_0$ is the biggest index $j < |\rho|$ such that $\rho(j) = 1$ (simillarly, if no such index exists, take $\tau$ instead of $\tau_{\rho \uh j_0}$). 

%    A node is said to be \emph{valid} if $\pi(\rho)$ is a condition.
     %A node $\rho$ is said to be \emph{valid} if the biggest index $i$ such that $\rho(i) = 0$ satisfies $\sigma_{\rho \uh i} \subseteq U$ (if no such index exists this condition is verified) and if the biggest index $j$ such that $\rho(i) = 1$ satisfies $\tau_{\rho \uh i} \subseteq V$ (if no such index exists this condition is verified). Being valid is a $\Delta_2^0(Y)$ condition.

     A node $\rho$ is said to be \emph{valid} if for every $\mu \in T$ such that $\mu \cdot 0 \preceq \rho$, $\sigma_{\mu} \subseteq U$ and for every $\mu \in T$ such that $\mu \cdot 1 \preceq \rho$, $\tau_{\mu} \subseteq V$. Being valid is a $\Delta^0_2(Y)$ condition.
\end{definition}

\begin{definition}
    Let $T$ a $(\sigma,\tau)$-tree 
    \begin{itemize}
        \item A node $(\sigma_{\rho}, \tau_{\rho})$ in $T$ \emph{satisfies} a pair of $\Sigma^0_1(Y)$-formulas $(\varphi_0(G_0, x), \varphi_1(G_1,x))$ if $\varphi_0(\sigma, |\rho| - \card \rho)$ and $\varphi_1(\tau, \card \rho)$ holds (where $\card \rho$ is the cardinal of $\rho$ seen as a set).
        \item $T$ \emph{satisfies} $(\varphi_0(G_0, x), \varphi_1(G_1,x))$ if every node of $T$ satisfies it.
    \end{itemize}
\end{definition}

\begin{lemma}\label[lemma]{lem:sads-wf-finite-tree-forcing}
Let $(\sigma,\tau)$ be a condition, $(\varphi_0(G_0, x), \varphi_1(G_1,x))$ be a pair of $\Sigma^0_1(Y)$ formulas, and $T$ be an $M$-finite maximal $(\sigma, \tau)$-tree satisfying $(\varphi_0(G_0, x), \varphi_1(G_1,x))$.
Then there is an extension $(\hat \sigma, \hat \tau) \leq (\sigma, \tau)$ and some~$p,q \in M \cup \{-1\}$ such that the following holds:
\begin{itemize}
    \item $(\hat \sigma, \hat \tau) \Vdash_0 \varphi_0(G_0, p)$ and $(\hat \sigma, \hat \tau) \Vdash_1 \varphi_1(G_1,q)$
    \item $(\hat \sigma, \hat \tau) \Vdash_0 \neg \varphi_0(G_1,p+1)$ or $(\hat \sigma, \hat \tau) \Vdash_1 \neg \varphi_1(G_1,q+1)$.
\end{itemize}
With the convention that  $(\hat \sigma, \hat \tau) \Vdash_i \varphi_i(G^i, -1)$ always holds.

\end{lemma}

\begin{proof}
If $T$ is empty, then $(\sigma,\tau) \nqvdash \varphi_0(G_0,0) \vee \varphi_1(G_1,0)$ as otherwise the corresponding split pair could be added in position $\epsilon$ in $T$, contradicting its maximality. Thus, by \Cref{lem:sads-forcing-question-sigma01}, there exists some extension $(\hat \sigma, \hat \tau) \leq (\sigma, \tau)$ and some side $i < 2$ such that $(\hat \sigma, \hat \tau) \Vdash_i \neg \varphi_i(G_i,0)$. Thus, we can assume that $T$ is not empty.


Consider a valid $\rho \in T$ having no valid descendant in $T$, this is possible since the set of valid element of $T$ is non empty because $(\sigma_\epsilon, \tau_\epsilon)$ is a vacuously valid, so we can take a maximal valid element of $T$ using $\mathsf{L}\Delta^0_2$ (which we have by $\BSig_2$).

Since $(\sigma_\rho, \tau_\rho)$ is a split pair, we have two cases:

Case 1: $\sigma_{\rho} \subseteq U$. In that case $\rho \cdot 0 \notin T$ by maximality of $\rho$. Then take $p = |\rho| - \card \rho$, $q = \card \rho - 1$ and $(\sigma', \tau') = (\sigma_{\rho}, \tau_{\rho \uh (i_0+1)})$ with $i_0$ the biggest index $i < |\rho|$ such that $\rho(i) = 1$ (if no such index exists we take $\tau$ instead of $\tau_{\rho \uh (i_0+1)}$), the pair $(\sigma', \tau')$ is a valid condition by definition of validity and the hypothesis $\sigma_{\rho} \subseteq U$. $(\sigma', \tau') \Vdash_0 \varphi_0(G_0, p)$ and $(\sigma', \tau') \Vdash_1 \varphi_1(G_1,q)$ because $T$ satisfies $(\varphi_0(G_0, x), \varphi_1(G_1,x))$. 

We have $(\sigma', \tau') \nqvdash \varphi_0(G_0,p + 1) \vee \varphi_1(G_1,q + 1)$ as otherwise the corresponding split pair could be added in position $\rho \cdot 0$ in $T$, contradicting its maximality. Hence, by \Cref{lem:sads-forcing-question-sigma01}, there exists some extension $(\hat \sigma, \hat \tau)$ and such that either $(\hat \sigma, \hat \tau) \Vdash_0 \neg \varphi_0(G_0,p + 1)$ or $(\hat \sigma, \hat \tau) \Vdash_1 \neg \varphi_1(G_1,q+1)$ holds.

%And $(\hat \sigma, \hat \tau) \Vdash_0 \neg \varphi_0(G_0,p + 1)$ by maximality of $T$ (Otherwise we could add $\rho \cdot 0$ to $T$ indexed by a split pair contradicting $(\hat \sigma, \hat \tau) \Vdash_0 \neg \varphi_0(G_0,p + 1)$).

Case 2: The case $\tau_{\rho} \subseteq V$ can be treated similarly. \\
\end{proof}

\end{document}
